\section{Introduction}

Let $N$ be a positive integer and $\Delta$ a subgroup with $\{\pm1\}\subseteq\Delta\subseteq(\Z/N\Z)^\times$.
The intermediate modular curve $X_\Delta(N)$ is the modular curve attached to
\[
 \Gamma_\Delta(N)=\left\{\begin{pmatrix}a&b\\c&d\end{pmatrix}\in\SL_2(\Z): c\equiv 0 \bmod N,\ a \bmod N\in\Delta\right\}.
\]
It lies between $X_1(N)=X_{\{\pm1\}}(N)$ and $X_0(N)=X_{(\Z/N\Z)^\times}(N)$, and it is defined over $\Q$.
We call $X_\Delta(N)$ \emph{strict} if both inclusions $\{\pm1\}\subsetneq\Delta\subsetneq(\Z/N\Z)^\times$
are strict. We follow the notation of Jeon, Kim and Schweizer \cite{JeonKimSchweizer2020}, Orli\'c
\cite{Orlic2025} and Derickx and Orli\'c \cite{DerickxOrlic2026}, and we name curves by their LMFDB labels
\cite{LMFDB,RSZB2022}.

For a curve $C$ over a field $k$, the \emph{$k$-gonality} $\gon_k C$ is the least degree of a nonconstant
morphism $C\to\PP^1$ defined over $k$. A \emph{point of degree $d$} on $C/\Q$ is a closed point whose residue
field has degree $d$ over $\Q$. The \emph{density degree set} $\delta(C/\Q)$ is the set of $d$ such that
$C$ has infinitely many points of degree $d$, and its least element $\min\delta(C/\Q)$ is the \emph{minimum
density degree} of $C$ \cite{KadetsVogt2025}. By Hilbert irreducibility, $C$ has infinitely many points of
degree $\gon_\Q C$, so $\min\delta(C/\Q)\le\gon_\Q C$.

\subsection{Previous work}
Points of low degree on modular curves have been studied intensively, with the degree going up one at a
time. For $X_0(N)$, the curves with infinitely many points of degree $2$ follow from Ogg's classification of
hyperelliptic curves \cite{Ogg1974} and Bars' classification of bielliptic curves \cite{Bars1999}; degree
$3$ is due to Jeon \cite{Jeon2021}; degree $4$ to Hwang and Jeon \cite{HwangJeon2024} and independently to
Derickx and Orli\'c \cite{DerickxOrlicX0}; and degree $5$ to Derickx, Hwang, Jeon and Orli\'c \cite{DHJO},
except for thirty levels that their methods did not reach. For intermediate curves, degree $2$ is due to
Jeon, Kim and Schweizer \cite{JeonKimSchweizer2020}, degree $3$ to Dalal \cite{Dalal2023}, and degree
$4$ to Derickx and Orli\'c \cite{DerickxOrlic2026}. In \cite{DerickxOrlic2026} the authors explain the obstacles to extending
their methods to degree $5$; Theorem~B below settles that case.

For gonality, the hyperelliptic and trigonal intermediate curves were found by Ishii and Momose
\cite{IshiiMomose1991} and Jeon and Kim \cite{JeonKim2007}; the list of Ishii and Momose misses the hyperelliptic
curve $X_\Delta(21)$ with $\Delta=\langle-1,8\rangle$, see \cite{JeonKim2007,JeonKimSchweizer2020}. Orli\'c \cite{Orlic2025} determined all
intermediate curves of $\Q$-gonality at most $5$, and the $\Q$-gonality of every strict curve with $N\le40$.
The $\Q$-gonality of $X_0(N)$ is known for $N<145$ by Najman and Orli\'c \cite{NajmanOrlic2024}, and that
of $X_1(N)$ for $N\le40$ by Derickx and van Hoeij \cite{DerickxVanHoeij2014}, who also gave upper bounds
for $N\le 250$.

\subsection{Results}
\begin{theoremC}\label{thm:C}
Let $X_\Delta(N)$ be a strict intermediate curve with $N\le 50$. Its $\Q$-gonality is given in
\cref{tab:C}. For $N\le 40$ these values are due to Orli\'c \cite{Orlic2025}, as are those of $6$ of the
$24$ curves with $41\le N\le 50$, which have gonality $4$ or $5$; two more have genus $3$ and $4$ and are
trigonal. The values for the other $16$ curves are new.
\end{theoremC}

By Orli\'c's classification of curves of $\Q$-gonality at most $5$ \cite{Orlic2025}, these $16$ curves have
$\Q$-gonality at least $6$; for the five of them with $\Q$-gonality $6$, the new part is the upper bound.

\begin{theoremD}\label{thm:D}
For $41\le N\le 50$, the $\Q$-gonality of $X_1(N)$ equals the upper bound of Derickx and van Hoeij
\cite{DerickxVanHoeij2014}:
\begin{center}
\begin{tabular}{l|cccccccccc}
$N$ & 41 & 42 & 43 & 44 & 45 & 46 & 47 & 48 & 49 & 50\\\hline
$\gon_\Q X_1(N)$ & 22 & 12 & 24 & 15 & 18 & 19 & 29 & 16 & 21 & 15
\end{tabular}
\end{center}
\end{theoremD}

Previously, Derickx and van Hoeij showed $\gon_\Q X_1(N)>8$ for $41\le N\le50$ and $\gon_\Q X_1(43)\ge17$
\cite[\S5]{DerickxVanHoeij2014}, and Najman and the author obtained lower bounds over finite fields,
$\gon_{\Fp}X_1(N)\ge10$ for $N\in\{41,42,44,45,46,48,50\}$ and $\ge19$ for $N=43$
\cite[Proposition~5.4 and Table~1]{NajmanVarivoda}.

\begin{theoremA}\label{thm:A}
For each of the thirty levels
\begin{align*}
 N\in\{&74, 77, 79, 82, 83, 85, 86, 88, 89, 92, 96, 98, 99, 100, 101,\\
 &103, 107, 111, 118, 123, 128, 131, 141, 142, 143, 145, 155, 159, 167, 191\}
\end{align*}
left open in \cite[Remark~1.11 and \S7]{DHJO}, the curve $X_0(N)$ has only finitely many points of degree $5$, and no
morphism of degree $5$ to $\PP^1$ over $\Q$. Together with \cite[Theorem~1.10]{DHJO}, this completes the
list of curves $X_0(N)$ with infinitely many points of degree $5$.
\end{theoremA}

The absence of a morphism of degree $5$ to $\PP^1$ was already known at $21$ of these levels \cite[\S7.2 and
Table~2]{DHJO}; it is new at $74,77,79,85,88,96,98,100,103$. The finiteness of the points of degree $5$ is new at all
thirty levels.

\begin{theoremB}\label{thm:B}
Exactly $43$ strict intermediate curves $X_\Delta(N)$ have infinitely many points of degree $5$. They are
listed in \cref{tab:positive}, and each of them has a morphism of degree $5$ to $\PP^1$ over $\Q$.
\end{theoremB}

\begin{corollaryE}\label{cor:E}
The strict intermediate curves of minimum density degree $5$ are exactly the four $\Q$-pentagonal curves of
Orli\'c \cite[Theorems~1.4--1.5]{Orlic2025}: $X_\Delta(31)$ with $\Delta=\langle-1,2\rangle$ and
$\Delta=\langle-1,5\rangle$, $X_\Delta(44)$ with $\Delta=\langle-1,5\rangle$, and $X_\Delta(125)$ with
$\Delta=\langle-1,4\rangle$.
\end{corollaryE}

This parallels the case of $X_0(N)$, where $X_0(109)$ is the only curve of minimum density degree $5$ \cite{DHJO}.

\subsection{Methods}
We have tried to keep the new mathematical input to a minimum, so that the paper can be checked by a human
reader. Wherever a published result or an inequality settles a case, we use it: Castelnuovo--Severi, the
classifications of curves of small gonality \cite{JeonKim2007,JeonKimSchweizer2020,Orlic2025}, and the
classifications of curves with infinitely many points of degree at most $4$
\cite{JeonKimSchweizer2020,Dalal2023,DerickxOrlic2026,DHJO}. Everything else is done by three methods,
each stated once:
\begin{enumerate}
\item[I.] \emph{Explicit pencils}: an explicit divisor $D$ with $\ell(D)\ge2$ gives $\gon_\Q\le\deg D$
 (\S\ref{sec:method-I}).
\item[II.] \emph{Finite-field exclusion}: at a prime of good reduction, a complete search over $\Fp$ shows
 that there is no function of degree at most $d$ (\S\ref{sec:method-II}). We use the gonality algorithm of
 Derickx and Terao \cite{DerickxTerao} and their package CurveArith. For the large curves $X_1(N)$ we keep
 their first phase (places and expansions of differentials) unchanged and replace their second phase by a
 search over a smaller family of divisors, run in compiled code (\S\ref{sec:phase2}). For $X_1(48)$ we
 search instead on the quotient by a diamond involution (\S\ref{sec:double-cover}).
\item[III.] \emph{Lifting arguments}: when a finite field allows more than $\Q$ does, a Mordell--Weil or
 local argument decides which classes over $\Fp$ come from $\Q$. This is needed for Theorems~A
 and~B and for one curve in Theorem~C.
\end{enumerate}

All computations were done in Magma~V2.29-4 \cite{Magma} and SageMath~10.7. The code, data and logs
that verify all computations in this paper can be found on
\begin{center}
 \url{https://github.com/marin-varivoda/intermediate-modular-curves};
\end{center}
\S\ref{sec:computations} describes how to rerun them.
